\chapter{总结与展望}
\section{研究总结}
本文以佐佐木希为主题，构建了覆盖五个职业阶段的分析框架，并从视觉材料与模拟评论两方面讨论跨媒介形象。阶段划分用于描述长期职业定位，情感指标用于比较短期反馈，二者共同形成纵向分析的基本结构。

在所设定的情境中，模拟评论在四个窗口呈现不同的情感构成，稳定更新期得分最高，调整期得分最低。由于结果来自预设样本，这一趋势只反映案例内部的计算关系。图片分析还表明，服饰、姿态与构图可以形成不同的视觉表达，但其与职业定位之间的关系仍需要同期文本支持。

\section{后续研究方向}
将本例发展为真实课题时，可先收集经核验的公开资料，明确时间边界与抽样方法，再开展独立编码与误差分析。不同平台与不同个案的比较，有助于检验阶段划分和情感指标的适用范围。

本例的时间线、事件与数值均为教学设定。附录\ref{app:checklist}提供复核清单，后续附录保留机械模型、带导数的运动方程与参数表，展示其他专业使用同一模板时的公式排版方式。
